\section{Formalism (Completed)}
\label{sec:formalism}
The proposed representation has three requirements. It must retain the
reparandum without asserting it as part of the speaker's final commitment,
identify the repair as accepted content and connect both candidates to the
structure they share. Meeting these requirements depends on SBN's treatment of
semantic roles, contexts and accessibility. Section~\ref{sec:sbn-brief}
introduces only the notation needed for this account; readers familiar with
SBN contexts and accessibility may begin with
Section~\ref{sec:paired-configuration}. Section~\ref{sec:formalism-coverage}
separates the representation's coverage from the narrower scope implemented
in the main experiment.

\subsection{SBN Requirements for the Representation}
\label{sec:sbn-brief}

SBN represents meaning as a sequence of concepts and their relations
\citep{bosSequenceNotationCatching2023}. Concepts introduce entities, events
or properties. Their labels record type and sense. In
Figure~\ref{fig:weather-rsr-comparison}(a), \sbnterm{city.n.01} introduces a
city concept and \sbnterm{Name "Dublin"} records its name. Roles connect concepts by
relative indices: \sbnterm{Location +1} points forward to the city, while
\sbnterm{Theme -1} points back to \sbnterm{weather.n.01}. Expressions such as
\sbnterm{EQU now} assign a value. The sequence therefore describes the request
without using explicit variables.

\begin{figure}[H]
\centering
\begin{minipage}[t]{0.53\textwidth}
\raggedright
\textbf{(a) No RSR}\par\smallskip
\parbox[t][3\baselineskip][t]{\linewidth}{\small\raggedright
\textit{``What is the weather like in Dublin?''}}
\begin{lstlisting}[style=sbn,
    basicstyle=\ttfamily\fontsize{7}{9}\selectfont,
    xleftmargin=0pt, xrightmargin=0pt, framexleftmargin=3pt,
    columns=fullflexible]
entity.n.01 EQU ?                                   % What
time.n.08 EQU now                                   % is
weather.n.01                                        % the weather
like.a.01 Co-Theme -3 Time -2 Theme -1 Location +1  % like 
city.n.01 Name "Dublin"                             % in Dublin?
\end{lstlisting}
\end{minipage}\hfill
\begin{minipage}[t]{0.42\textwidth}
\raggedright
\textbf{(b) Replacement self-repair}\par\smallskip
\parbox[t][3\baselineskip][t]{\linewidth}{\small\raggedright
\textit{``What is the weather like in Berlin, in Dublin?''}}
\begin{lstlisting}[style=sbn,
    basicstyle=\ttfamily\fontsize{7}{9}\selectfont,
    xleftmargin=0pt, xrightmargin=0pt, framexleftmargin=3pt,
    columns=fullflexible, escapeinside={(*@}{@*)}]
entity.n.01 EQU ?                       % What
time.n.08 EQU now                       % is
weather.n.01                            % the weather
like.a.01 Co-Theme -3 Time -2 Theme -1  % like 
    (*@\colorbox{blue!10}{\textcolor{blue!60!black}{\bfseries CORRECTION <1}}@*)
city.n.01 Name "Berlin" LocationOf -1   % in Berlin
    (*@\colorbox{blue!10}{\textcolor{blue!60!black}{\bfseries CONJUNCTION <2}}@*)                    % I mean, 
city.n.01 Name "Dublin" LocationOf -2   % in Dublin?
\end{lstlisting}
\end{minipage}
\caption{Fluent and repair-aware SBNs for the weather example. The highlighted
separators introduce the withdrawn Berlin material and the accepted Dublin
material.}
\label{fig:weather-rsr-comparison}
\end{figure}

The fluent request occupies one \emph{context}, which groups concepts with the
conditions that describe them. A \emph{separator} starts a new context when
part of the meaning requires separate interpretation. Its \emph{connector}
points to the context to which the new one relates. SBN has no closing
separator, so subsequent concepts remain in the new context until another
separator occurs or the sequence ends. For example, \sbnterm{NEGATION} places
the buying event and car in a context whose content is denied, while John
remains in the enclosing context
\citep{bosVariablefreeDiscourseRepresentation2021}:

\begin{lstlisting}[style=sbn]
male.n.02 Name "John"                  % John
    NEGATION <1                        % did not
time.n.08 TPR now
buy.v.01 Agent -2 Time -1 Theme +1     % buy
car.n.01                               % a car
\end{lstlisting}

The \sbnterm{Agent -2} link reaches out from the negated context to John. A
link in the opposite direction cannot directly reach the car introduced only
inside that context. These \emph{accessibility} restrictions determine which
concepts can be linked across context boundaries. They are central to RSR
because the reparandum must connect to shared content without becoming part of
the accepted proposition.

\subsection{The Paired RSR Configuration}
\label{sec:paired-configuration}

I propose a representation that pairs the existing separators
\sbnterm{CORRECTION} and \sbnterm{CONJUNCTION}; it introduces no new
separator.\footnote{A fuller
specification and additional examples are available at
\url{https://hongxuzhou.github.io/lct_master_project/}.} In
Figure~\ref{fig:weather-rsr-comparison}(b), \sbnterm{CORRECTION} introduces the
context containing Berlin, while \sbnterm{CONJUNCTION} introduces the context
containing Dublin. The two connectors point to the initial context containing
the shared request. Berlin and Dublin therefore occupy separate contexts with
the same parent.

The second separator is necessary because SBN does not close a context
automatically. Without it, Dublin and the continuation would remain inside
the withdrawn context. \sbnterm{CONJUNCTION} starts the accepted context and
merges its content with the shared parent
\citep{bosSequenceNotationCatching2023}. Following the treatment of corrected
content as denied in SDRT
\citep{lascaridesAgreementDisputesCommitments2009}, I give the
\sbnterm{CORRECTION} context an implicit negative interpretation. Berlin is
withdrawn; Dublin is accepted.

Both cities remain connected to the weather request. Berlin carries
\sbnterm{LocationOf -1} from the withdrawn context, and Dublin carries
\sbnterm{LocationOf -2} from the accepted context. These links identify the
role being revised. Berlin's location relation remains inside the withdrawn
content, so it does not entail that Berlin is the requested location. Dublin's
relation contributes to the accepted proposition. The representation thus
preserves the two candidates and records the change in their discourse status.

The paired configuration is realised through three constructions: the direct,
anchor and movement constructions. The direct construction requires each
candidate to carry a suitable link to shared content. A bare-subject RSR such
as \textit{``Josh, I mean, Mary
bought a cake''} lacks such a link on the subject concept itself. The
\emph{anchor construction} places a neutral \sbnterm{entity.n.01} in the shared
context and links Josh and Mary to it with \sbnterm{EQU}. The buying event
keeps one agent; Josh is withdrawn as its identification and Mary is accepted.
The PMB already uses neutral entity concepts for underspecified participants
\citep{abzianidzeParallelMeaningBank2017}. Appendix~\ref{app:devices-anchor}
gives the full construction.

A different problem occurs when the repair concept points forward to an
argument. In \textit{``Tom opened, no, broke the window''}, a copy of the
verb inside the withdrawn context cannot reach forward through both separators
to \textit{the window}. The \emph{movement construction} places the window before
the paired RSR contexts, allowing both verbs to link back to it. This preserves the
shared argument but departs from surface word order. Direct construction is
tried first. Where it fails, referential nouns can use the anchor construction,
while properties and events use the movement construction when their arguments
can be moved safely.
Appendix~\ref{app:devices} gives the listings and selection rules.

The treatment also supports later reference to accessible withdrawn
referents. Proper names can be interpreted outside the scope of negation
through presupposition
\citep[Section~7]{bosVariablefreeDiscourseRepresentation2021}. Berlin and
Dublin can therefore serve as antecedents for \textit{these two cities}, while
Berlin's location role remains withdrawn. Later reference does not restore
Berlin as the accepted location. Appendix~\ref{app:later-reference} gives the
complete \sbnterm{ANA} links. The same account does not make every item inside
a withdrawn context accessible; the challenge set is restricted to referents
for which the formal account licenses access.

Following \citet{maschlerMetalanguagingDiscourseMarkers1994},
\textit{I mean} is treated as a metalinguistic marker when it functions as an
editing expression. It is excluded from the propositional content under repair
and recorded in a \sbnterm{\%} comment on the \sbnterm{CONJUNCTION} line, as
illustrated in Figure~\ref{fig:weather-rsr-comparison}(b).

\subsection{Coverage, Implementation and Remaining Validation}
\label{sec:formalism-coverage}

The configuration can represent replacement of content expressed by SBN
concepts or roles, provided that the two candidates can connect to a shared
parent without violating accessibility. This includes replacements inside an
existing context such as negation. The corpus used in the main experiment has
a narrower implemented scope: it transforms aligned noun, verb and adjective
concepts. Role replacement is part of the representational account but is not
included in the present training data.

The paired configuration does not cover replacement of an SBN separator.
Under the current SBN rules, separators organise relations between contexts
rather than contribute content within them, so the
\sbnterm{CORRECTION}--\sbnterm{CONJUNCTION} structure cannot embed competing
separators as the reparandum and repair. Representing such cases would require
modifying SBN itself and falls outside the scope of the proposed thesis. This
limitation applies to changes in modality, such as \textit{must} to
\textit{might}, and changes in discourse relation, such as \textit{when} to
\textit{even though}, because they revise the structures that organise the
contexts. Tense and aspect replacements can also require several coordinated
concept changes and are excluded from the current corpus.

Timing sets a further boundary. The configuration covers an immediate RSR,
where the repair follows before the speaker commits intervening material.
Consider \textit{``I bought three apples from Lidl yesterday---actually,
four.''} Starting the withdrawn context at \textit{three} would also place the
location and time inside it. Moving that committed span before the context
would depart too far from the utterance order. Delayed self-repairs of this kind,
including corrections made after a completed clause, fall outside the scope of
the proposed thesis. They could be forced into the current notation only by
moving intervening committed material across contexts, at the cost of a
substantial departure from word order and SBN's preference for alignment with
the natural-language input.

The paired configuration and the direct, anchor and movement constructions
have been proposed and implemented. The completed corpus pass in
Section~\ref{sec:data} initially tests their structural applicability. The
thesis will complete the answer to the first research
question stated in Section~\ref{sec:introduction} by validating the retained corpus, examining where the constructions
succeed or fail and consolidating the resulting coverage and limitation
account. The parsing experiment then tests whether models can learn the
distinctions that the representation makes explicit.